\documentclass[11pt,a4paper]{article}
\usepackage[margin=2.6cm]{geometry}
\usepackage{amsmath,amssymb,amsthm}
\usepackage[T1]{fontenc}
\usepackage{lmodern}
\usepackage{microtype}
\usepackage{enumitem}
\setlength{\parindent}{0pt}
\setlength{\parskip}{0.65em}
\title{\textbf{MATH6310 Coursework 1}}
\author{Sai Vanam}
\date{}
\begin{document}
\maketitle
\vspace{-1.5em}

\section*{Exercise 0.1 --- Discrete metric \hfill \normalsize (10 marks)}
\textbf{Question.} Prove that the discrete metric on a set $S$ is a metric.

\textbf{Solution.}
The discrete metric states that, for any set $S$, we define
\[
 d_{\mathrm{discrete}}(s,t)=
 \begin{cases}
 0, & \text{if }s=t,\\
 1, & \text{if }s\ne t.
 \end{cases}
\]
To prove that this is a metric, it has to satisfy three conditions for every $x,y,z\in S$:
\begin{align*}
\text{(MS1)}\qquad &d(x,y)=0\iff x=y,\\
\text{(MS2)}\qquad &d(x,y)=d(y,x)\quad\text{(symmetry)},\\
\text{(MS3)}\qquad &d(x,z)\le d(x,y)+d(y,z)\quad\text{(triangle inequality)}.
\end{align*}
Also, $d$ only takes the values $0$ and $1$, so its values are always non-negative and finite.

\textbf{(MS1).} This is satisfied because, by the definition of the discrete metric, $d(x,y)=0$ exactly when $x=y$. If $x\ne y$, the distance is $1$, not $0$.

\textbf{(MS2).} To prove symmetry, consider any $x,y\in S$.
If $x=y$, then $d(x,y)=0=d(y,x)$.
If $x\ne y$, then $y\ne x$ as well, so $d(x,y)=1=d(y,x)$.
Therefore, in both cases,
\[
 d(x,y)=d(y,x).
\]

\textbf{(MS3).} To prove the triangle inequality, consider any $x,y,z\in S$.

If $x=z$, then $d(x,z)=0$, and since distances are non-negative,
\[
 0=d(x,z)\le d(x,y)+d(y,z).
\]

If $x\ne z$, then $d(x,z)=1$. In this case, $x\ne y$ or $y\ne z$ (or both): otherwise $x=y$ and $y=z$ would imply $x=z$, which is a contradiction. Therefore, at least one of $d(x,y)$ and $d(y,z)$ equals $1$, so
\[
 d(x,y)+d(y,z)\ge 1=d(x,z).
\]
Thus the triangle inequality holds in both cases.

Since (MS1), (MS2) and (MS3) are satisfied, the discrete metric is a metric on any set $S$. \(\square\)

\section*{Exercise 0.4 --- Counting metric \hfill \normalsize (10 marks)}
\textbf{Question.} Verify the Triangle Inequality for the counting metric.

\textbf{Solution.}
Let $S$ be a finite set. For any subsets $A,B\subseteq S$, the counting metric is defined by
\[
 d(A,B)=\#(A\triangle B),
\]
where $\#(A\triangle B)$ means the number of elements in the \emph{symmetric difference}
\[
 A\triangle B=(A\setminus B)\cup(B\setminus A).
\]
In other words, $A\triangle B$ contains all elements that belong to exactly one of $A$ and $B$.

We need to show that, for any $A,B,C\subseteq S$,
\[
 d(A,C)\le d(A,B)+d(B,C).
\]
First, we claim that
\[
 A\triangle C\subseteq (A\triangle B)\cup(B\triangle C).
\]
To see this, let $x\in A\triangle C$. Then $x$ belongs to exactly one of $A$ and $C$.

If $x\in A$ and $x\notin C$, there are two possibilities. If $x\notin B$, then $x\in A\triangle B$. If $x\in B$, then $x\in B\triangle C$.

Similarly, if $x\in C$ and $x\notin A$, then either $x\in B$, in which case $x\in A\triangle B$, or $x\notin B$, in which case $x\in B\triangle C$.

Hence every element of $A\triangle C$ lies in at least one of $A\triangle B$ and $B\triangle C$, proving the claimed inclusion.

Because $S$ is finite, taking the number of elements of both sides gives
\begin{align*}
\#(A\triangle C)
&\le \#\bigl((A\triangle B)\cup(B\triangle C)\bigr)\\
&\le \#(A\triangle B)+\#(B\triangle C).
\end{align*}
The second inequality holds because the number of elements in a union is at most the sum of the numbers of elements in the two sets (elements in both are not counted twice in the union).

Using the definition of the counting metric, this becomes
\[
 \boxed{d(A,C)\le d(A,B)+d(B,C).}
\]
Therefore the counting metric satisfies the Triangle Inequality (MS3). \(\square\)

\end{document}
